\chapter{K\"all\'en--Lehmann Spectral Representation}
\label{ch:kallen-lehmann}

\section{Spectral Data of a Scalar Field}

Let \(\Hilb\) carry a strongly continuous unitary representation of spacetime
translations with commuting self-adjoint generators \(P^\mu\). Let
\[
  E(\Delta)
\]
be the joint projection-valued spectral measure of \(P^\mu\), where
\(\Delta\subset\mathbb R^D\) is Borel. Assume the spectrum condition
\[
  \operatorname{Spec}(P)\subset\overline V_+
  =
  \{p\in\mathbb R^D:p^0\ge0,\ p^2\le0\},
\]
with \(p^2=\eta_{\mu\nu}p^\mu p^\nu\). Let \(\Omega\) be a translation
invariant vacuum vector.

Let \(\widehat\phi\) be a scalar operator-valued distribution. If
\[
  \langle\Omega|\widehat\phi(x)|\Omega\rangle
\]
is nonzero, replace \(\widehat\phi\) by
\(\widehat\phi-\langle\Omega|\widehat\phi|\Omega\rangle\). This removes the
vacuum atom from the two-point function and keeps the connected spectral data.
Point-field notation below is distributional notation; after smearing, the
matrix elements are ordinary Hilbert-space pairings.

The Wightman two-point distribution is
\[
  W_2(x-y)
  =
  \langle\Omega|
  \widehat\phi(x)\widehat\phi(y)
  |\Omega\rangle.
\]
Translation covariance gives
\[
  W_2(x-y)
  =
  \int_{\overline V_+}\ee^{\ii p\cdot(x-y)}\,\dd\nu_\phi(p),
\]
where the positive measure \(\nu_\phi\) is
\[
  \dd\nu_\phi(p)
  =
  \langle\Omega|
  \widehat\phi(0)E(\dd p)\widehat\phi(0)
  |\Omega\rangle.
\]
Positivity is immediate from the spectral theorem: for a Borel set
\(\Delta\),
\[
  \nu_\phi(\Delta)
  =
  \|E(\Delta)\widehat\phi(0)\Omega\|^2
  \ge0.
\]

\section{Invariant-Mass Decomposition}

For a scalar field and invariant vacuum, \(\nu_\phi\) is invariant under the
proper orthochronous Lorentz group. The Lorentz-invariant parameter on the
forward cone is the invariant mass squared
\[
  \mu^2=-p^2\ge0.
\]
Let
\[
  \pi:\overline V_+\to[0,\infty),
  \qquad
  \pi(p)=-p^2,
\]
and define the positive Borel measure
\[
  \rho(B)=\nu_\phi(\pi^{-1}(B)).
\]
This is the pushforward of \(\nu_\phi\) by the invariant-mass map.  By
disintegration of positive measures, \(\nu_\phi\) is an integral over
conditional measures supported on the level sets \(\pi^{-1}(\mu^2)\).  Lorentz
invariance fixes each nonzero conditional measure to be proportional to the
standard invariant measure on the positive mass shell.  With the normalization
absorbed into \(\dd\rho(\mu^2)\), this gives
\[
  \dd\nu_\phi(p)
  =
  \int_{[0,\infty)}
  \dd\rho(\mu^2)\,
  \frac{\dd^Dp}{(2\pi)^{D-1}}\,
  \theta(p^0)\delta(p^2+\mu^2).
\]
Here \(d\rho\) is a positive Borel measure on \([0,\infty)\). When it is
absolutely continuous one writes \(\dd\rho(\mu^2)=\rho(\mu^2)\dd\mu^2\);
when it has atoms, the atomic weights are part of the same measure.

Define the positive-frequency scalar two-point function of mass \(\mu\) by
\[
  \Delta_+(x;\mu^2)
  =
  \int\frac{\dd^Dp}{(2\pi)^{D-1}}\,
  \theta(p^0)\delta(p^2+\mu^2)\,
  \ee^{\ii p\cdot x}.
\]
Equivalently,
\[
  \Delta_+(x;\mu^2)
  =
  \int
  \frac{\dd^{D-1}\vec p}{(2\pi)^{D-1}}\,
  \frac{
    \ee^{\ii\vec p\cdot\vec x-\ii\sqrt{\vec p^{\,2}+\mu^2}\,x^0}
  }{2\sqrt{\vec p^{\,2}+\mu^2}} .
\]
The Wightman two-point function is then
\[
  W_2(x)
  =
  \int_{[0,\infty)}
  \dd\rho(\mu^2)\,
  \Delta_+(x;\mu^2).
\]
This is the K\"all\'en--Lehmann representation for the Wightman two-point
function.

\section{Time-Ordered and Euclidean Forms}

The time-ordered two-point function is
\[
  G_T(x)
  =
  \langle\Omega|
  \operatorname{T}\{
    \widehat\phi(x)\widehat\phi(0)
  \}
  |\Omega\rangle.
\]
Using the definition of time ordering,
\[
  G_T(x)
  =
  \theta(x^0)W_2(x)+\theta(-x^0)W_2(-x).
\]
Therefore
\[
  G_T(x)
  =
  \int_{[0,\infty)}
  \dd\rho(\mu^2)\,
  \Delta_F(x;\mu^2),
\]
where
\[
  \Delta_F(x;\mu^2)
  =
  -\ii
  \int\frac{\dd^Dk}{(2\pi)^D}
  \frac{\ee^{\ii k\cdot x}}{k^2+\mu^2-\ii\epsilon}.
\]
Equivalently,
\[
  G_T(x)
  =
  -\ii
  \int\frac{\dd^Dk}{(2\pi)^D}\,
  \ee^{\ii k\cdot x}
  \int_{[0,\infty)}
  \frac{\dd\rho(\mu^2)}{k^2+\mu^2-\ii\epsilon}.
\]

The Euclidean two-point function has the same spectral measure:
\[
  G_E(x_E)
  =
  \int_{[0,\infty)}
  \dd\rho(\mu^2)\,
  \Delta_E(x_E;\mu^2),
\]
with
\[
  \Delta_E(x_E;\mu^2)
  =
  \int\frac{\dd^Dk_E}{(2\pi)^D}
  \frac{\ee^{\ii k_E\cdot x_E}}{k_E^2+\mu^2}.
\]
The spectral measure \(d\rho\) is thus common to the Wightman, time-ordered,
and Euclidean two-point functions. The different Green functions are different
boundary values or analytic continuations of the same positive spectral data.

\section{Stieltjes Positivity of the Euclidean Propagator}

The Euclidean momentum-space two-point function contains the same information
in a form that is often the most robust for constructive and lattice
arguments.  Set
\[
  Q^2=k_E^2\ge0 .
\]
After Fourier transform,
\begin{equation}
  \widetilde G_E(Q^2)
  =
  \int_{[0,\infty)}
  \frac{\dd\rho(s)}{Q^2+s}.
  \label{eq:euclidean-stieltjes-propagator}
\end{equation}
Thus \(\widetilde G_E\) is a Stieltjes transform of a positive measure.
Consequently, wherever the derivatives exist,
\[
  (-1)^r
  \frac{\dd^r}{\dd (Q^2)^r}\widetilde G_E(Q^2)
  =
  r!\int_{[0,\infty)}
  \frac{\dd\rho(s)}{(Q^2+s)^{r+1}}
  \ge0,
  \qquad r=0,1,2,\ldots .
\]
These inequalities are not perturbative statements.  They are Hilbert-space
positivity written in Euclidean momentum variables.

The measure can be recovered from the discontinuity of the analytic function
\[
  z\longmapsto
  \int_{[0,\infty)}\frac{\dd\rho(s)}{z+s}
\]
across its cut on the negative real \(z\)-axis.  If the absolutely continuous
part has density \(\rho_{\rm ac}(s)\), then
\[
  \rho_{\rm ac}(s)
  =
  -\frac{1}{2\pi\ii}
  \operatorname{Disc}_{z=-s}\widetilde G_E(z),
  \qquad s>0,
\]
with atoms appearing as simple poles at \(z=-s_0\) of positive residue.  This
is the Euclidean form of the statement that stable particles are isolated
positive atoms of the spectral measure and multiparticle states are continuous
spectral support.

\section{Canonical Normalization and Spectral Weight}

The measure \(d\rho\) depends on the normalization of the local field.  When
\(\widehat\phi\) is a canonically normalized scalar field, the equal-time
commutation relation supplies an additional normalization condition.  Assume
that, on a common invariant domain,
\[
  [\widehat\phi(t,\vec x),\partial_t\widehat\phi(t,\vec y)]
  =
  \ii\,\delta^{D-1}(\vec x-\vec y)\mathbf 1 .
\]
The commutator two-point distribution is
\[
  \langle\Omega|[\widehat\phi(x),\widehat\phi(0)]|\Omega\rangle
  =
  \int_{[0,\infty)}\dd\rho(\mu^2)\,
  \bigl(\Delta_+(x;\mu^2)-\Delta_+(-x;\mu^2)\bigr).
\]
For each fixed \(\mu^2\), the equal-time derivative of the expression in
parentheses gives the canonical delta distribution.  Therefore the canonical
commutator fixes the total spectral weight:
\begin{equation}
  \int_{[0,\infty)}\dd\rho(\mu^2)=1.
  \label{eq:kallen-lehmann-canonical-sum-rule}
\end{equation}
This is a statement about the chosen field normalization, not about a
normalization-independent observable.

If the same field has a one-particle atom \(Z\delta_{m^2}\), positivity of
\(d\rho\) and \eqref{eq:kallen-lehmann-canonical-sum-rule} imply
\[
  0\le Z\le1.
\]
The remainder \(1-Z\) is the total spectral weight carried by all continuum
states and by any other atoms.  For a noncanonical interpolating field, a
composite operator, or a renormalized field in a different convention, the
same K\"all\'en--Lehmann representation holds, but the right-hand side of
\eqref{eq:kallen-lehmann-canonical-sum-rule} is replaced by the corresponding
equal-time commutator normalization.

\section{One-Particle Atoms}

Assume the spectrum contains an isolated scalar one-particle mass \(m>0\), and
assume the local field \(\widehat\phi\) has nonzero projection onto that
one-particle subspace. With the invariant normalization from the one-particle
construction,
\[
  \langle\vec p\,{}'|\vec p\rangle
  =
  (2\pi)^{D-1}2\omega_{\vec p}\,
  \delta^{D-1}(\vec p\,{}'-\vec p),
  \qquad
  \omega_{\vec p}=\sqrt{\vec p^{\,2}+m^2},
\]
define \(Z\ge0\) by
\[
  \langle\vec p|\widehat\phi(0)|\Omega\rangle=Z^{1/2}
\]
for scalar one-particle states. Lorentz covariance makes the right-hand side
independent of \(\vec p\) in this normalization. The one-particle contribution
to the spectral measure is the atom
\[
  Z\,\delta_{m^2},
\]
where \(\delta_{m^2}\) is the unit point mass at \(\mu^2=m^2\). Thus
\[
  \dd\rho(\mu^2)
  =
  Z\,\delta_{m^2}(\dd\mu^2)
  +
  \dd\rho_{\mathrm c}(\mu^2),
\]
where \(d\rho_{\mathrm c}\) is the remaining positive measure.

\begin{center}
\begin{tikzpicture}[>=Stealth,scale=0.95]
  \node (v) at (0,0) {$\widehat\phi(f)\Omega$};
  \node[draw,rounded corners=2pt,inner sep=4pt,blue!70!black] (vac) at (3.0,1.25)
    {$E(\{0\})\widehat\phi(f)\Omega$};
  \node[draw,rounded corners=2pt,inner sep=4pt,green!45!black] (one) at (3.25,0)
    {$E(\mu^2=m^2)\widehat\phi(f)\Omega$};
  \node[draw,rounded corners=2pt,inner sep=4pt,orange!80!black] (cont) at (3.35,-1.25)
    {$E_{\mathrm c}\widehat\phi(f)\Omega$};
  \draw[->,thick] (v) -- (vac);
  \draw[->,thick] (v) -- (one);
  \draw[->,thick] (v) -- (cont);
  \node[blue!70!black] at (6.3,1.25) {vacuum atom};
  \node[green!45!black] at (6.6,0) {one-particle atom};
  \node[orange!80!black] at (6.6,-1.25) {continuum spectral part};
\end{tikzpicture}
\end{center}

After subtracting the vacuum expectation value of \(\widehat\phi\), the vacuum
projection in the first line is absent. The one-particle atom remains precisely
when the field has nonzero overlap with the isolated mass shell.

Under the same assumptions,
\[
  G_T(x)
  =
  Z\,\Delta_F(x;m^2)
  +
  \int_{[0,\infty)}
  \dd\rho_{\mathrm c}(\mu^2)\,
  \Delta_F(x;\mu^2).
\]
If the theory has a single stable particle of mass \(m\), no lighter
particles, and the continuum part is generated by states containing at least
two such particles, then
\[
  \operatorname{supp}d\rho_{\mathrm c}\subset[4m^2,\infty).
\]
When \(d\rho_{\mathrm c}\) has a density \(\sigma(\mu^2)\) in that region,
\[
  G_T(x)
  =
  Z\,\Delta_F(x;m^2)
  +
  \int_{4m^2}^{\infty}
  \dd\mu^2\,\sigma(\mu^2)\Delta_F(x;\mu^2).
\]

\begin{center}
\begin{tikzpicture}[>=Stealth,scale=0.95]
  \draw[->,line width=0.7pt] (0,0) -- (6.4,0) node[right] {$s=\mu^2$};
  \draw[->,line width=0.7pt] (0,0) -- (0,3.05) node[above] {$\dd\rho(s)$};
  \draw (1.65,0.08) -- (1.65,-0.08) node[below] {$m^2$};
  \draw (3.35,0.08) -- (3.35,-0.08) node[below] {$4m^2$};
  \draw[red!75!black,very thick,->] (1.65,0) -- (1.65,2.45)
    node[above] {$Z\,\delta_{m^2}$};
  \draw[blue!65!black,thick]
    (3.35,0) .. controls (3.7,0.65) and (4.15,1.20) .. (4.75,1.12)
    .. controls (5.35,1.04) and (5.75,0.48) .. (6.05,0.30);
  \node[blue!65!black,align=center] at (4.85,2.25)
    {continuum support\\under the stated threshold assumption};
\end{tikzpicture}
\end{center}

The number \(Z\) is attached to the chosen local field. Rescaling the field
rescales \(Z\). Once the field normalization is fixed, \(Z\) is a definite
matrix element and the one-particle pole residue of its two-point function.

\section{Role in the Construction}

The K\"all\'en--Lehmann representation is nonperturbative. It states that
Hilbert-space positivity and the spectrum condition force the scalar
two-point function to be a positive superposition of free massive two-point
functions. Isolated atoms of the measure identify stable one-particle mass
shells seen by the local field. Continuous spectral support records the
remaining invariant-mass spectrum.

The next step in the construction is the perturbative computation of Green
functions. The construction of scattering states and the S-matrix comes later,
after the spectral one-particle data have been separated from the local
two-point functions.
